\documentclass[11pt,a4paper]{article}

\usepackage[a4paper,margin=2.2cm,top=2.5cm,bottom=2.5cm]{geometry}
\usepackage{fontspec}
\usepackage{amsmath,amssymb}
\usepackage{enumitem}
\usepackage{array}
\usepackage{microtype}
\usepackage{titling}

\setmainfont{Latin Modern Roman}[Ligatures=TeX]

\setlength{\parindent}{0pt}
\setlength{\parskip}{6pt}

\begin{document}

% ==============================================================================
% INSTITUTIONAL EXAM HEADER
% ==============================================================================
\begin{center}
    {\large\bfseries Sylhet Engineering College, Sylhet}\\[3pt]
    {\normalsize\bfseries Department of Computer Science and Engineering}\\[3pt]
    {\normalsize\bfseries Term Test-01}\\[6pt]
\end{center}

\vspace{2pt}
\noindent
\textbf{Course:} CSE 0541 1143 (Discrete Mathematics) \hfill \textbf{Marks:} 10\\
\rule{\linewidth}{0.8pt}

\vspace{10pt}

% ==============================================================================
% QUESTIONS
% ==============================================================================
\begin{enumerate}[label=\textbf{\arabic*.}, leftmargin=1.5em, itemsep=18pt]

    % Question 1
    \item 
    \begin{enumerate}[label=\textbf{(\alph*)}, leftmargin=1.8em, itemsep=10pt]
        \item Define a \textbf{proposition} and determine whether each of the following is a proposition. Give reasons. \hfill \textbf{[03]}
        \begin{enumerate}[label=\roman*), leftmargin=2em, itemsep=4pt]
            \item $x + 5 = 10$
            \item Dhaka is the capital of Bangladesh.
        \end{enumerate}

        \item Construct the truth table for the compound proposition: \hfill \textbf{[02]}
        \[
        (p \to q) \wedge \neg q
        \]
    \end{enumerate}

    % Question 2
    \item Let $A = \{1, 2, 3, 4\}$ and a relation $R$ on $A$ by:
    \[
    R = \{(a, b) \in A \times A \mid a \text{ divides } b\}.
    \]
    \begin{enumerate}[label=\textbf{(\alph*)}, leftmargin=1.8em, itemsep=10pt]
        \item Write the relation $R$ explicitly as a set of ordered pairs. \hfill \textbf{[02]}

        \item Determine whether $R$ is \textbf{reflexive}, \textbf{symmetric}, \textbf{antisymmetric}, and \textbf{transitive}. Provide a brief justification for each property. \hfill \textbf{[03]}
    \end{enumerate}

\end{enumerate}

\vfill
\begin{center}
    \rule{0.3\linewidth}{0.4pt}\\[2pt]
    {\footnotesize\itshape --- End of Question Paper ---}
\end{center}

\end{document}
